\newglossaryentry{api}{
    name=API,
    description={Application Programming Interface, interface de programmation permettant à plusieurs applications ou composants logiciels de communiquer entre eux}
}

\newglossaryentry{bdmr}{
    name=BDMR,
    description={Base de données des modèles raster, dépôt historique utilisé au SHOM pour archiver et versionner les productions cartographiques raster exploitées par la chaîne Ediacara}
}

\newglossaryentry{ci}{
    name=CI,
    description={Continuous Integration, ou intégration continue, pratique consistant à intégrer régulièrement les modifications de code dans un dépôt partagé et à automatiser les vérifications associées}
}

\newglossaryentry{cd}{
    name=CD,
    description={Continuous Delivery ou Continuous Deployment selon le contexte, démarche visant à automatiser la mise à disposition ou le déploiement des nouvelles versions d'une application}
}

\newglossaryentry{crud}{
    name=CRUD,
    description={Create, Read, Update, Delete, ensemble des quatre opérations de base appliquées à la gestion de données persistées}
}

\newglossaryentry{devops}{
    name=DEVOPS,
    description={Approche rapprochant développement et exploitation afin d'améliorer l'automatisation, la qualité, la fréquence des livraisons et la collaboration autour du cycle de vie logiciel}
}

\newglossaryentry{dto}{
    name=DTO,
    description={Data Transfer Object, objet utilisé pour transporter des données entre couches ou services sans exposer directement le modèle de persistance}
}

\newglossaryentry{enc}{
    name=ENC,
    description={Electronic Navigational Chart, carte électronique de navigation normalisée destinée aux systèmes de navigation maritime}
}

\newglossaryentry{fcu}{
    name=FCU,
    description={Fond de carte unifié, mode de production cartographique du SHOM visant à automatiser davantage la génération des cartes papier à partir des données disponibles dans ses bases}
}

\newglossaryentry{gdal}{
    name=GDAL,
    description={Geospatial Data Abstraction Library, bibliothèque et suite d'outils de traitement et de conversion de données géospatiales raster et vectorielles}
}

\newglossaryentry{geotiff}{
    name=GEOTIFF,
    description={Format TIFF enrichi d'informations de géoréférencement permettant d'associer une image raster à un système de coordonnées géographiques}
}

\newglossaryentry{gitops}{
    name=GITOPS,
    description={Pratique d'exploitation où l'état souhaité d'une infrastructure ou d'applications est décrit dans Git puis appliqué automatiquement par des outils de déploiement}
}

\newglossaryentry{ia}{
    name=IA,
    description={Intelligence artificielle, ensemble de méthodes permettant à un système informatique de réaliser des tâches d'analyse, de recommandation ou de décision à partir de règles ou de données}
}

\newglossaryentry{ipf}{
    name=IPF,
    description={Innovation Project Framework, cadre pédagogique du projet de fin d'année dans lequel a été développé Lynx Immo avec un client réel et une équipe d'étudiants de M1 et M2}
}

\newglossaryentry{jwt}{
    name=JWT,
    description={JSON Web Token, format compact de jeton signé utilisé notamment pour transporter des informations d'authentification et d'autorisation}
}

\newglossaryentry{kpi}{
    name=KPI,
    description={Key Performance Indicator, indicateur clé de performance utilisé pour mesurer l'avancement, la qualité, l'usage ou l'efficacité d'un projet ou d'un produit}
}

\newglossaryentry{mvp}{
    name=MVP,
    description={Minimum Viable Product, première version d'un produit regroupant les fonctionnalités essentielles permettant de démontrer sa valeur et de recueillir des retours avant de poursuivre son enrichissement}
}

\newglossaryentry{orm}{
    name=ORM,
    description={Object-Relational Mapping, mécanisme permettant de faire correspondre des objets applicatifs à des tables d'une base de données relationnelle}
}

\newglossaryentry{owasp}{
    name=OWASP,
    description={Open Worldwide Application Security Project, fondation produisant des ressources et bonnes pratiques de référence pour la sécurité des applications}
}

\newglossaryentry{pdl}{
    name=PDL,
    description={Plan de Développement Logiciel, document de gouvernance décrivant l'organisation, la méthodologie, les rôles, le planning, les outils, la qualité et les principaux choix techniques d'un projet logiciel}
}

\newglossaryentry{poc}{
    name=POC,
    description={Proof of Concept, preuve de concept destinée à valider la faisabilité ou la pertinence d'une solution avant une industrialisation éventuelle}
}

\newglossaryentry{pss}{
    name=PSS,
    description={Project Scope Statement, document de cadrage précisant le périmètre, les objectifs, les livrables, les jalons, les hypothèses, les contraintes et les critères d'acceptation d'un projet}
}

\newglossaryentry{raci}{
    name=RACI,
    description={Responsible, Accountable, Consulted, Informed, matrice de responsabilités permettant d'identifier pour une activité qui la réalise, qui en porte la responsabilité finale, qui doit être consulté et qui doit être informé}
}

\newglossaryentry{rest}{
    name=REST,
    description={Representational State Transfer, style d'architecture utilisé pour concevoir des services web autour de ressources manipulées au moyen des standards du protocole HTTP}
}

\newglossaryentry{rgpd}{
    name=RGPD,
    description={Règlement général sur la protection des données, cadre réglementaire européen relatif à la protection des données à caractère personnel}
}

\newglossaryentry{rls}{
    name=RLS,
    description={Row Level Security, mécanisme de sécurité de PostgreSQL permettant de contrôler l'accès aux lignes d'une table selon des politiques définies}
}

\newglossaryentry{rsx}{
    name=RSX,
    description={Radiosignaux, famille d'ouvrages nautiques du SHOM regroupant différentes informations radioélectriques utiles à la navigation et à la sécurité maritime}
}

\newglossaryentry{s100}{
    name=S-100,
    description={Cadre de l'Organisation hydrographique internationale définissant une architecture commune pour les nouvelles générations de produits et services hydrographiques numériques}
}

\newglossaryentry{s123}{
    name=S-123,
    description={Spécification de produit fondée sur S-100 destinée à représenter les services radio maritimes}
}

\newglossaryentry{shom}{
    name=SHOM,
    description={Service hydrographique et océanographique de la Marine, établissement public français de référence pour l'hydrographie, l'océanographie et l'information géographique maritime et littorale}
}

\newglossaryentry{smdsm}{
    name=SMDSM,
    description={Système mondial de détresse et de sécurité en mer, dispositif international de radiocommunication dédié à la sécurité maritime}
}

\newglossaryentry{svn}{
    name=SVN,
    description={Apache Subversion, système centralisé de gestion de versions utilisé notamment pour historiser des fichiers et des arborescences de production}
}

\newglossaryentry{ui}{
    name=UI,
    description={User Interface, interface utilisateur regroupant les éléments visuels et interactifs par lesquels une personne utilise une application}
}

\newglossaryentry{uml}{
    name=UML,
    description={Unified Modeling Language, langage de modélisation normalisé utilisé pour représenter la structure et le comportement d'un système logiciel}
}

\newglossaryentry{ux}{
    name=UX,
    description={User Experience, expérience utilisateur correspondant à la manière dont une personne perçoit et utilise un produit ou un service numérique}
}

\newglossaryentry{xsd}{
    name=XSD,
    description={XML Schema Definition, langage décrivant la structure, les types et les contraintes applicables à des documents XML}
}